% 06_discussion.tex
\section{Discussion}
\label{sec:discussion}

\subsection{Root cause of the Pearson $r$ gap}
\label{sec:disc_precision}

Our reproduction recovers $r = 0.864$ between $f^+_{cp}$ and $H_p$
across all 23{,}391 fragments, falling $0.106$ below the paper's target
of $r \approx 0.97$.
We attribute this gap entirely to the floating-point precision of pLDDT
values in AlphaFold-DB~v4.

In the paper's dataset --- most likely an earlier AlphaFold-DB version
where pLDDT appears to have been stored or rounded to integer values ---
a typical 300-residue protein has $\lesssim\!20$ distinct pLDDT levels,
each populated by many residues forming a substantial horizontal plateau
in $N_{cc}(\mathrm{pLDDT})$.
In AlphaFold-DB~v4, the same protein can carry over 100 distinct
floating-point levels, so each level is typically occupied by only 1--3
residues.
This eliminates the broad plateaus that make $N_{cc}$ peaks persistent.

The effect on $H_p$ is additive: more distinct persistence values in
$\mathcal{P}_D$ spread the probability mass function more evenly,
increasing $H_p$ by a roughly constant offset ($+0.14$ to $+0.33$ across
our three prototypes, Table~\ref{tab:prototypes}).
Both $f^+_{cp}$ and $H_p$ are shifted upward by comparable amounts, so
the cloud of points in Figure~\ref{fig:scatter_corr} translates
diagonally without rotating.
A diagonal translation preserves the slope but increases the residual
variance around the regression line, reducing $r$ from the paper's $0.97$
to our $0.864$.

\subsection{Root cause of the PLM candidate gap}
\label{sec:disc_plm}

The 183-vs-86 discrepancy in fragmentation candidates follows from the
same mechanism, but acts on a different statistic.
PLM counts how many $N_{cc}$ peaks are \emph{persistent} at
$t_\nu = n \cdot t_p$, measured in connected-component-count units.
A peak requires a substantial local maximum in $N_{cc}$ that collapses
slowly as pLDDT decreases.

With integer-valued pLDDT, 20--30 residues may share the same level,
creating a wide plateau where $N_{cc}$ stays elevated for a long
filtration interval.
With float pLDDT and one-residue-per-step insertion, those residues are
spread across many individual levels, so the $N_{cc}$ curve has many
small local fluctuations.
Each distinct float level can create a local maximum in $N_{cc}$ when
a new singleton is inserted before merging into an existing component.
These fine-grained peaks often have persistence exceeding $t_\nu$, which
is why we observe more candidates ($183$) than the paper ($86$) rather
than fewer.

This explains why the gap is approximately multiplicative (${\approx}2\times$):
float-level granularity creates additional persistent peaks beyond those
detected by the paper's integer-level algorithm.
The 34 proteins surviving the strictest threshold $t_p = 0.030$ constitute
a robust core whose domain boundaries produce exceptionally large $N_{cc}$
peaks regardless of filtration granularity.

\subsection{Biological interpretation of the arity enrichment}
\label{sec:disc_arity}

The 5.92-fold enrichment of fragmented proteins in the arity bin
$(a_{25}=7, a_{75}=15)$ is statistically significant
($p = 1.61 \times 10^{-3}$) and biologically interpretable.
The arity range $[7, 15]$ at the 25th--75th percentile level corresponds
to proteins that are neither maximally disordered nor maximally compact:

\begin{itemize}
  \item A low-arity $a_{25} \approx 7$ means that a quarter of the
  protein's residues have fewer than 7 $\mathrm{C}_\alpha$ neighbours
  within 10\,\AA{}.
  This is characteristic of flexible linker regions with limited
  tertiary contacts.

  \item A moderate $a_{75} \approx 15$ means that three-quarters of
  residues have at most 15 neighbours, which is typical of well-packed
  globular domains but below the $a_{75} \approx 20$ seen in
  tightly-packed $\alpha/\beta$ proteins.
\end{itemize}

Together, this signature is exactly what one expects from a multi-domain
protein: compact, independently folded domains contribute the upper
quartile of arities, while disordered linkers contribute the lower
quartile.
Purely disordered proteins (like A0A0G2L439, $a_{75}=7$) have low arity
throughout, and single-domain globular proteins (like P15121, $a_{75}=23$)
have uniformly high arity.
Multi-domain fragmented proteins inhabit the intermediate zone where
packing heterogeneity is highest.

Crucially, this enrichment is derived from the 3D coordinate geometry
of the AlphaFold structure, entirely independently of the pLDDT
sequence analysis.
The fact that two orthogonal structural descriptors --- topological
pLDDT fragmentation and $\mathrm{C}_\alpha$ packing arity --- converge
on the same set of proteins provides mutual validation for both methods.

\subsection{Limitations}
\label{sec:disc_limits}

\paragraph{Single organism.}
All results are computed on the \textit{H.~sapiens} proteome only.
The paper demonstrates similar correlations across multiple model
organisms; whether the arity enrichment pattern generalises to other
proteomes remains untested.

\paragraph{No experimental disorder ground truth.}
We do not compare our fragmented candidates against DisProt or MobiDB
experimental annotations of intrinsic disorder.
Such a comparison would establish whether PLM$\,\geq 3$ proteins are
enriched in annotated IDRs or in multi-domain entries in the literature,
providing a biological validation beyond the internal cross-correlation.
This is out of scope for the current course project.

\paragraph{Data version.}
AlphaFold-DB~v6 was released after this project began; however, we
subsequently re-ran the full pipeline on the v6 tarball as a verification
step (Section~\ref{sec:disc_v6}).
All proteome-scale statistics change by less than $0.5\%$ between v4 and v6,
confirming that the pLDDT representation is identical between the two releases
and that the float-precision gap is not version-specific.

\subsection{AlphaFold-DB v6 replication}
\label{sec:disc_v6}

To test whether the float-precision explanation is specific to
AlphaFold-DB~v4 or persists in the current release, we re-ran the full
pipeline on the AlphaFold-DB~v6 \textit{H.~sapiens} proteome tarball
(5.1\,GB; downloaded May~2026).
Table~\ref{tab:v4v6} summarises the comparison.

\begin{table}[h]
  \caption{Proteome-scale statistics: AlphaFold-DB v4 vs.\ v6.}
  \label{tab:v4v6}
  \centering
  \small
  \begin{tabular}{@{}lrrr@{}}
    \toprule
    Metric & v4 & v6 & $\Delta$ \\
    \midrule
    $N$ proteins           & 23{,}391 & 23{,}586 & $+195$   \\
    Pearson $r$            & 0.864    & 0.862    & $-0.002$ \\
    Median $f^+_{cp}$      & 0.242    & 0.239    & $-0.003$ \\
    Median $H_p$           & 0.362    & 0.359    & $-0.003$ \\
    Candidates$^\dagger$   & 183      & 185      & $+2$     \\
    \bottomrule
  \end{tabular}\\[3pt]
  \footnotesize $^\dagger$ $n \geq 200$, $H_p \geq 0.25$,
  $\mathrm{PLM}(t_p{=}0.025) \geq 3$.
\end{table}

The v6 release adds 196 proteins and removes 150 relative to v4 (UniProt
curation updates), but all five statistics are statistically
indistinguishable.
The Pearson $r$ changes by $-0.002$, the fragmentation candidate count by
$+2$, and the median $H_p$ and $f^+_{cp}$ by less than $0.4\%$ in absolute
terms.
This confirms that our discrepancy with the paper is not a version artifact:
both v4 and v6 store pLDDT as floating-point values and yield the same
elevated $H_p$ profile relative to the paper's effectively integer-quantised
data.

\subsection{Directions for future work}
\label{sec:disc_future}

The most direct extension would be to apply the same pipeline to additional
model organisms
(\textit{M.\ musculus}, \textit{S.\ cerevisiae}) to assess the
universality of the arity enrichment signal.
A comparison of our 183 fragmented candidates against UniProt domain
annotations and known multi-domain proteins would provide external
biological validation.
On the algorithmic side, the sublevel-set filtration framework could be
generalised to other scalar functions on sequence-embedded graphs,
such as predicted solvent accessibility or AlphaFold's PAE matrix;
persistence of the resulting $N_{cc}$ curves could yield
complementary fragmentation descriptors.

